% =====================================================================
%  Resolução detalhada — Sistemas Elétricos de Potência
%  Questões 51, 24 e 45, em texto corrido (A4): o que nos slides
%  (beamer/resolucao_grei.tex) é dito oralmente vem aqui por extenso,
%  com as deduções e as contas passo a passo.
%  Compilar a partir desta pasta (documento/):
%      latexmk -pdf resolucao_detalhada.tex     (pdflatex + biber)
%  A bibliografia é a mesma dos slides: ../beamer/referencias.bib
% =====================================================================
\documentclass[a4paper,11pt]{article}

\usepackage[utf8]{inputenc}
\usepackage[T1]{fontenc}
\usepackage[brazil]{babel}
\usepackage{lmodern}
\usepackage{microtype}
\usepackage[a4paper,top=2.5cm,bottom=2.5cm,left=2.5cm,right=2.5cm,
            headheight=14pt]{geometry}
\usepackage{amsmath,amssymb}
\usepackage{icomma}            % vírgula decimal em modo matemático: 0,5
\usepackage{booktabs}
\usepackage{array}
\usepackage[shortlabels]{enumitem}
\usepackage{graphicx}
\usepackage{pifont}            % \ding: marcas de certo/errado
\usepackage{float}             % [H]: figura exatamente onde está no texto
\usepackage{caption}
\usepackage[american,cuteinductors]{circuitikz}
\usetikzlibrary{arrows.meta,calc}
\usepackage{pgfplots}          % gráfico da potência instantânea (Q45)
\pgfplotsset{compat=1.18}
\providecommand{\sen}{\operatorname{sen}}   % "sen", como no enunciado da Q45
\usepackage[most]{tcolorbox}
\usepackage{titlesec}
\usepackage{fancyhdr}
\usepackage{csquotes}
\usepackage[backend=biber,style=numeric,sorting=none]{biblatex}
\addbibresource{../beamer/referencias.bib}
% deixa a URL quebrar em qualquer letra/número (sem isso ela vaza da margem)
\setcounter{biburllcpenalty}{7000}
\setcounter{biburlucpenalty}{8000}
\setcounter{biburlnumpenalty}{9000}
\usepackage[hidelinks]{hyperref}

% ---------------------------------------------------------------------
%  Cores: as do template GREI (TemplateBeamerUFC/libs/ufc_format.cls)
%  + os neutros das figuras dos slides
% ---------------------------------------------------------------------
\definecolor{blue_theme}{RGB}{47, 93, 75}     % verde GREI (estrutural)
\definecolor{red}{RGB}{215, 0, 0}             % alerta/resposta
\definecolor{green_ufc}{RGB}{81, 203, 238}    % azul claro (auxiliar)
\definecolor{Ink}{RGB}{34, 38, 43}
\definecolor{Mute}{RGB}{122, 130, 140}
\colorlet{ZonaA}{green_ufc}
\colorlet{ZonaB}{blue_theme}
\colorlet{CorP}{blue_theme}   % Q45: potência ativa (média)
\colorlet{CorS}{red}          % Q45: potência aparente (amplitude) / reativa

% ---------------------------------------------------------------------
%  Página: títulos, cabeçalho, numeração
% ---------------------------------------------------------------------
\titleformat{\section}{\Large\bfseries\color{blue_theme}}{\thesection}{0.8em}{}
\titleformat{\subsection}{\large\bfseries\color{blue_theme}}{\thesubsection}{0.8em}{}
\titleformat{\subsubsection}{\normalsize\bfseries\color{Ink}}{\thesubsubsection}{0.8em}{}
\setcounter{secnumdepth}{3}
\setcounter{tocdepth}{2}
\numberwithin{equation}{section}
\numberwithin{figure}{section}
\numberwithin{table}{section}

\pagestyle{fancy}
\fancyhf{}
\fancyhead[L]{\small\color{Mute}Resolução detalhada --- Sistemas Elétricos de Potência}
\fancyhead[R]{\small\color{Mute}GREI/UFC}
\fancyfoot[C]{\small\thepage}
\renewcommand{\headrulewidth}{0.4pt}

\setlength{\parindent}{0pt}
\setlength{\parskip}{0.55em}
\setlist{itemsep=0.2em,topsep=0.3em}
\captionsetup{font=small,labelfont={bf,color=blue_theme}}

% ---------------------------------------------------------------------
%  Caixas
% ---------------------------------------------------------------------
% fórmula/conceito-chave
\newtcolorbox{conceito}[1][]{colback=blue_theme!4,colframe=blue_theme,
  boxrule=0.6pt,arc=2pt,left=8pt,right=8pt,top=4pt,bottom=4pt,
  coltitle=white,colbacktitle=blue_theme,fonttitle=\bfseries\small,#1}
% resposta
\newtcolorbox{resposta}[1][]{colback=red!4,colframe=red,boxrule=0.8pt,
  arc=2pt,left=8pt,right=8pt,top=5pt,bottom=5pt,
  coltitle=white,colbacktitle=red,fonttitle=\bfseries\small,
  title=Resposta,#1}
% observação/interpretação
\newtcolorbox{observacao}[1][]{colback=green_ufc!8,colframe=green_ufc!70!black,
  boxrule=0.6pt,arc=2pt,left=8pt,right=8pt,top=4pt,bottom=4pt,
  coltitle=white,colbacktitle=green_ufc!70!black,fonttitle=\bfseries\small,#1}

% rótulo de passo (como nos slides)
\newcommand{\passo}[1]{{\color{blue_theme}\bfseries #1}\enspace}
% marcas de certo/errado (tabela de alternativas)
\newcommand{\ok}{\textcolor{blue_theme}{\ding{51}}}
\newcommand{\nok}{\textcolor{red}{\ding{55}}}

% ---------------------------------------------------------------------
%  Símbolos para diagramas unifilares (os mesmos dos slides)
% ---------------------------------------------------------------------
\tikzset{
  barra/.style={line width=2.2pt,Ink},
  fio/.style={line width=0.8pt,Ink},
  gerador/.pic={
    \draw[fio,fill=white] (0,0) circle (0.36);
    \draw[fio] (-0.2,0) sin (-0.1,0.12) cos (0,0) sin (0.1,-0.12) cos (0.2,0);
  },
  trafo/.pic={
    \draw[fio] (-0.62,0) -- (-0.44,0) (0.44,0) -- (0.62,0);
    \fill[white] (-0.17,0) circle (0.27) (0.17,0) circle (0.27);
    \draw[fio] (-0.17,0) circle (0.27) (0.17,0) circle (0.27);
  },
  rot/.style={font=\footnotesize,Ink},
  numbarra/.style={circle,draw=Ink,line width=0.6pt,inner sep=1.2pt,
                   font=\scriptsize,fill=white},
}
\ctikzset{bipoles/length=1.05cm}

% ---------------------------------------------------------------------
%  Questão 51: unifilar (o mesmo desenho dos slides)
%  #1 = escala do tikzpicture; os raios dos pics não acompanham o scale,
%  então os rótulos G e T ficam a raio/escala abaixo do centro.
% ---------------------------------------------------------------------
\newcommand{\circuitoQcinquentaum}[1]{%
  \pgfmathsetmacro{\dG}{0.36/#1}%
  \pgfmathsetmacro{\dT}{0.30/#1}%
  \draw[fio] (0.36,2) -- (1.2,2);   \pic at (0,2) {gerador};
  \node[rot,anchor=north] at (0,2-\dG) {$G_1$};
  \draw[fio] (9.6,2) -- (10.44,2);  \pic at (10.8,2) {gerador};
  \node[rot,anchor=north] at (10.8,2-\dG) {$G_2$};
  \foreach \x in {1.2,3.6,7.2,9.6} \draw[barra] (\x,0.55) -- (\x,3.45);
  \node[rot,above] at (1.2,3.45) {20 kV};
  \node[rot,above] at (3.6,3.45) {250 kV};
  \node[rot,above] at (7.2,3.45) {250 kV};
  \node[rot,above] at (9.6,3.45) {25 kV};
  \foreach \y/\n in {3.0/1,1.0/2}{
    \draw[fio] (1.2,\y) -- (1.78,\y) (3.02,\y) -- (3.6,\y);
    \pic at (2.4,\y) {trafo};
    \node[rot,anchor=north,inner sep=1pt] at (2.4,\y-\dT) {$T_{\n}$};}
  \foreach \y/\n in {3.0/3,1.0/4}{
    \draw[fio] (7.2,\y) -- (7.78,\y) (9.02,\y) -- (9.6,\y);
    \pic at (8.4,\y) {trafo};
    \node[rot,anchor=north,inner sep=1pt] at (8.4,\y-\dT) {$T_{\n}$};}
  \draw[fio] (3.6,2.4) -- node[rot,above] {$j125\,\Omega$} (7.2,2.4);
  \draw[fio] (3.6,1.6) -- node[rot,below] {$j125\,\Omega$} (7.2,1.6);
}
\newcommand{\unifilarQcinquentaum}{%
\begin{tikzpicture}[scale=0.85]\circuitoQcinquentaum{0.85}\end{tikzpicture}}

% zonas de tensão base (estáticas: todas de uma vez)
% \zonaQcinquentaum{x ini}{x fim}{cor}{nome}{V_b em kV}{elementos}
\newcommand{\zonaQcinquentaum}[6]{%
  \fill[#3!13] (#1,-0.05) rectangle (#2,4.05);
  \node[anchor=south,font=\small,text=#3!75!black,inner sep=2pt]
    at ({(#1+#2)/2},4.05) {\textbf{#4}: $V_b = #5$~kV};
  \node[anchor=north,align=center,font=\footnotesize,text=Ink,inner sep=2pt]
    at ({(#1+#2)/2},-0.05) {#6};}
% fronteira tracejada, interrompida onde ficam os rótulos T
\newcommand{\fronteiraQcinquentaum}[1]{%
  \draw[Ink!75,line width=0.7pt,dash pattern=on 3pt off 2pt]
    (#1,4.05) -- (#1,3-\dT) (#1,2.45-\dT) -- (#1,1-\dT);}
\newcommand{\zonasQcinquentaum}{%
\begin{tikzpicture}[scale=0.95]
  \zonaQcinquentaum{-1.35}{2.4}{ZonaA}{Zona 1}{20}{$G_1$ e lado de\\20 kV de $T_1$, $T_2$}
  \zonaQcinquentaum{2.4}{8.4}{ZonaB}{Zona 2}{250}{lado de 250 kV de $T_1$ a $T_4$\\e as duas linhas}
  \zonaQcinquentaum{8.4}{12.15}{ZonaA}{Zona 3}{25}{lado de 25 kV de\\$T_3$, $T_4$ e $G_2$}
  \circuitoQcinquentaum{0.95}
  \fronteiraQcinquentaum{2.4}
  \fronteiraQcinquentaum{8.4}
\end{tikzpicture}}

% circuito equivalente em pu; [#1] = escala; #2 = 0 completo, 1 sem a linha 2
\newcommand{\equivQcinquentaum}[2][0.9]{%
\begin{circuitikz}[scale=#1, transform shape, line width=0.7pt]
  \draw (0,0) to[sV, l=$E_1\angle\delta$] (0,2.4)
        to[L, l=$j0{,}1$] (2,2.4);
  \draw (10,0) to[sV, l_=$E_2\angle 0^\circ$] (10,2.4)
        to[L, l_=$j0{,}1$] (8,2.4);
  \foreach \x in {2,4,6,8} \draw (\x,1.7) -- (\x,3.1);
  \draw (2,3.1) to[L, l=$j0{,}2$] (4,3.1);
  \draw (2,1.7) to[L, l_=$j0{,}2$] (4,1.7);
  \draw (4,3.1) to[L, l=$j0{,}2$] (6,3.1);
  \ifnum#2=0
    \draw (4,1.7) to[L, l_=$j0{,}2$] (6,1.7);
  \else
    \draw[Mute!50, dashed] (4,1.7) -- (6,1.7);
    \draw[red, line width=1.4pt] (4.75,1.45) -- (5.25,1.95) (4.75,1.95) -- (5.25,1.45);
    \node[font=\footnotesize, red] at (5,1.05) {aberta};
  \fi
  \draw (6,3.1) to[L, l=$j0{,}2$] (8,3.1);
  \draw (6,1.7) to[L, l_=$j0{,}2$] (8,1.7);
  \draw (0,0) -- (10,0);
  \foreach \x/\t in {1/G_1, 3/T_1\parallel T_2, 5/L_1\parallel L_2, 7/T_3\parallel T_4, 9/G_2}
    \node[font=\footnotesize, blue_theme] at (\x,4.05) {$\t$};
  \foreach \x/\t in {2/20 kV, 4/250 kV, 6/250 kV, 8/25 kV}
    \node[font=\scriptsize, Mute, fill=white, inner sep=1pt] at (\x,0.55) {\t};
\end{circuitikz}}

% duas fontes ligadas por uma reatância (equação potência-ângulo)
\newcommand{\duasfontes}{%
\begin{circuitikz}[scale=0.95, transform shape, line width=0.7pt]
  \draw (0,0) to[sV, l=$E_1\angle\delta$] (0,1.8)
        to[L, l=$jX$] (3.6,1.8)
        to[sV, l=$E_2\angle 0^\circ$] (3.6,0) -- (0,0);
  \draw[red,line width=1pt,-{Latex[length=2mm]}] (1.2,1.2) -- (2.4,1.2)
    node[midway,below,font=\small] {$I$};
\end{circuitikz}}

% curva potência-ângulo genérica (x em graus)
\newcommand{\curvaPdelta}{%
\begin{tikzpicture}[x=0.035cm, y=1.5cm, >={Latex[length=2mm]}]
  \draw[->,Ink] (0,0) -- (200,0) node[right] {$\delta$};
  \draw[->,Ink] (0,0) -- (0,1.35) node[above] {$P$};
  \draw[Mute,dashed] (0,1) -- (90,1) -- (90,0);
  \draw[blue_theme,line width=1.1pt,domain=0:180,samples=61] plot (\x,{sin(\x)});
  \node[above,font=\small] at (90,1) {$P_{\max} = E_1E_2/X$};
  \node[below,font=\small] at (0,0) {$0$};
  \node[below,font=\small] at (90,0) {$90^\circ$};
  \node[below,font=\small] at (180,0) {$180^\circ$};
  \node[font=\footnotesize,Mute,align=center] at (45,0.35) {região\\estável};
  \node[font=\footnotesize,Mute,align=center] at (135,0.35) {região\\instável};
\end{tikzpicture}}

% curvas antes/depois da perda da linha, com delta = 30 graus
\newcommand{\curvasQcinquentaum}{%
\begin{tikzpicture}[x=0.045cm, y=1.45cm, >={Latex[length=2mm]}]
  \draw[->,Ink] (0,0) -- (195,0) node[right] {$\delta$};
  \draw[->,Ink] (0,0) -- (0,2.3) node[above] {$P$ (pu)};
  \draw[blue_theme,line width=1.1pt,domain=0:180,samples=91] plot (\x,{2*sin(\x)});
  \draw[red,line width=1.1pt,domain=0:180,samples=91] plot (\x,{sin(\x)/0.6});
  \draw[Mute,dashed] (0,1) -- (30,1) (0,0.8333) -- (30,0.8333) (30,0) -- (30,1);
  \draw[Mute,dashed] (0,2) -- (90,2) (0,1.6667) -- (90,1.6667) (90,0) -- (90,2);
  \fill[blue_theme] (30,1) circle (1.8pt);
  \fill[red] (30,0.8333) circle (1.8pt);
  \node[left,font=\footnotesize] at (0,2) {$2{,}0$};
  \node[left,font=\footnotesize] at (0,1.6667) {$1{,}667$};
  \node[left,font=\footnotesize] at (0,1.06) {$1{,}0$};
  \node[left,font=\footnotesize] at (0,0.78) {$0{,}833$};
  \node[below,font=\footnotesize] at (30,0) {$30^\circ$};
  \node[below,font=\footnotesize] at (90,0) {$90^\circ$};
  \node[below,font=\footnotesize] at (180,0) {$180^\circ$};
  \node[anchor=west,font=\small,align=left] at (200,1.3)
    {\textcolor{blue_theme}{\rule[0.5ex]{1.4em}{1.1pt}}\ antes: $P = 2{,}0\,\sin\delta$\\[0.3em]
     \textcolor{red}{\rule[0.5ex]{1.4em}{1.1pt}}\ depois: $P = 1{,}667\,\sin\delta$};
\end{tikzpicture}}

% linha k-m só com reatância
\newcommand{\linhaKM}{%
\begin{tikzpicture}[scale=1.0]
  \draw[barra] (0,0.3) -- (0,1.3);
  \draw[barra] (3.6,0.3) -- (3.6,1.3);
  \draw[fio] (0,0.8) -- node[rot,above] {$jx$} (3.6,0.8);
  \draw[red,line width=0.9pt,-{Latex[length=2mm]}] (1.1,0.55) -- (2.5,0.55)
    node[midway,below,font=\footnotesize] {$I_{km}$};
  \node[rot,below] at (0,0.3) {$V_k = |V_k|\angle\theta_k$};
  \node[rot,below] at (3.6,0.3) {$V_m = |V_m|\angle\theta_m$};
  \node[rot,above] at (0,1.3) {barra $k$};
  \node[rot,above] at (3.6,1.3) {barra $m$};
\end{tikzpicture}}

% Questão 24: unifilar; [#1] escala; #2 = 0 enunciado, 1 com a solução
\newcommand{\unifilarQvintequatro}[2][0.95]{%
\begin{tikzpicture}[scale=#1]
  \draw[fio] (0.36,1) -- (1.2,1);   \pic at (0,1) {gerador};
  \draw[fio] (7.6,1) -- (8.44,1);   \pic at (8.8,1) {gerador};
  \foreach \x in {1.2,4.4,7.6} \draw[barra] (\x,0.35) -- (\x,1.65);
  \foreach \x/\n in {1.2/1,4.4/2,7.6/3} \node[numbarra] at (\x,2.0) {\n};
  \draw[fio] (1.2,1) -- node[rot,above] {$j0{,}2$} (4.4,1);
  \draw[fio] (4.4,1) -- node[rot,above] {$j0{,}2$} (7.6,1);
  \draw[fio,-{Latex[length=2.2mm]}] (4.4,0.6) -- (4.85,0.6) -- (4.85,-0.15);
  \node[rot,anchor=north west] at (7.7,0.4) {$S_{G3}=0{,}5$};
  \ifnum#2=1
    \draw[red,line width=1pt,-{Latex[length=2mm]}] (2.0,0.62) -- (3.6,0.62);
    \node[font=\scriptsize,red,below] at (2.8,0.6) {$S_{12}=0{,}5$};
    \draw[red,line width=1pt,-{Latex[length=2mm]}] (6.8,0.62) -- (5.3,0.62);
    \node[font=\scriptsize,red,below] at (6.1,0.6) {$S_{32}=0{,}5$};
    \node[font=\scriptsize,red,anchor=north] at (4.85,-0.2) {$S_{c}=1{,}0-j0{,}1$};
    \node[font=\scriptsize,blue_theme,anchor=south] at (1.2,2.25) {$1\angle 0^\circ$};
    \node[font=\scriptsize,blue_theme,anchor=south] at (4.4,2.25) {$1{,}005\angle{-5{,}71^\circ}$};
    \node[font=\scriptsize,blue_theme,anchor=south] at (7.6,2.25) {$1\angle 0^\circ$};
  \fi
\end{tikzpicture}}

% Questão 24: diagrama fasorial em escala
\newcommand{\fasorialQvintequatro}{%
\begin{tikzpicture}[>={Latex[length=2.4mm]}]
  \draw[->,green_ufc!80!black,line width=1.1pt] (0,0) -- (10,0)
    node[above left,font=\small] {$V_3 = V_1 = 1\angle 0^\circ$};
  \draw[->,red,line width=1.1pt] (0,0.25) -- node[above,font=\small] {$I_{32} = 0{,}5\angle 0^\circ$} (5,0.25);
  \draw[->,Mute,line width=0.9pt] (10,0) -- (10,-1)
    node[right,font=\small,pos=0.6] {$-jx\,I_{32} = -j0{,}1$};
  \draw[->,blue_theme,line width=1.1pt] (0,0) -- (10,-1)
    node[below,font=\small,pos=0.55,yshift=-2pt] {$V_2 = 1{,}0 - j0{,}1$};
  \draw[Mute] (7,0) arc (0:-5.71:7);
  \node[font=\footnotesize,Mute,anchor=west] at (7.05,-0.33) {$\theta_2$};
\end{tikzpicture}}

% ---------------------------------------------------------------------
%  Questão 45: curva de potência instantânea (pgfplots)
%  p(t) = 500 - 1000 cos(2wt - 60°), w = 100*pi rad/s (f = 50 Hz);
%  em graus, 2wt = 36000 t.
%  \graficoQquarentacinco[opções do axis]{0 = enunciado | 1 = com P e S}
% ---------------------------------------------------------------------
\pgfplotsset{
  q45/.style={
    width=7cm, height=5.5cm,
    xmin=0, xmax=0.02, ymin=-600, ymax=1600,
    xtick={0,0.002,0.004,0.006,0.008,0.01,0.012,0.014,0.016,0.018,0.02},
    ytick={-600,-400,...,1600},
    scaled x ticks=false,
    grid=major, major grid style={densely dotted,Ink!40},
    xlabel={tempo em segundos}, ylabel={watt},
    tick label style={font=\tiny,Ink}, label style={font=\scriptsize,Ink},
    xticklabel style={/pgf/number format/.cd,fixed,precision=3,use comma,/tikz/.cd},
    yticklabel style={/pgf/number format/.cd,use comma,1000 sep={}},
    axis line style={Ink}, tick style={Ink!70},
  },
  % rótulos do tempo inclinados (slides, onde a figura é estreita)
  q45 girado/.style={xticklabel style={rotate=45,anchor=north east,inner sep=1pt}},
}
\newcommand{\graficoQquarentacinco}[2][]{%
\begin{tikzpicture}
\begin{axis}[q45,#1]
  \addplot[Ink,line width=1pt,domain=0:0.02,samples=241]
    {500 - 1000*cos(36000*x - 60)};
  \ifnum#2=1
    \draw[Mute,line width=0.5pt] (axis cs:0,1500) -- (axis cs:0.02,1500)
                                 (axis cs:0,-500) -- (axis cs:0.02,-500);
    \addplot[CorP,line width=1pt,dashed,domain=0:0.02] {500};
    \draw[CorS,line width=1pt,{Latex[length=1.8mm]}-{Latex[length=1.8mm]}]
      (axis cs:0.00667,500) -- (axis cs:0.00667,1500)
      node[midway,right,font=\scriptsize] {$S$};
    \draw[CorS,line width=1pt,{Latex[length=1.8mm]}-{Latex[length=1.8mm]}]
      (axis cs:0.01167,500) -- (axis cs:0.01167,-500)
      node[midway,left,font=\scriptsize] {$S$};
    \node[CorP,font=\scriptsize,fill=white,inner sep=1pt,anchor=west]
      at (axis cs:0.0128,500) {$P$};
    \node[font=\scriptsize,Ink,anchor=north west,inner sep=1.5pt]
      at (axis cs:0.0085,1500) {$p_{\max}\approx 1500$};
    \node[font=\scriptsize,Ink,anchor=south east,inner sep=1pt]
      at (axis cs:0.0198,-500) {$p_{\min}\approx -500$};
  \fi
\end{axis}
\end{tikzpicture}}
% triângulo de potências
\newcommand{\trianguloPotencias}{%
\begin{tikzpicture}[scale=0.95, >={Latex[length=2.2mm]}]
  \draw[->,CorP,line width=1.1pt] (0,0) -- (3.2,0)
    node[midway,below,font=\small] {$P = S\cos\theta$};
  \draw[->,CorS,line width=1.1pt] (3.2,0) -- (3.2,2.4)
    node[midway,right,font=\small] {$Q = S\sen\theta$};
  \draw[->,Ink,line width=1.1pt] (0,0) -- (3.2,2.4)
    node[midway,above left,font=\small] {$S = VI$};
  \draw[Mute] (0.9,0) arc (0:36.87:0.9);
  \node[font=\small] at (1.2,0.32) {$\theta$};
\end{tikzpicture}}

% decomposição p(t) = P(1 - cos 2wt) - Q sen 2wt, com P = 500 W e Q = 866 var
% (só no documento)
\newcommand{\decomposicaoQquarentacinco}{%
\begin{tikzpicture}
\begin{axis}[q45, width=11cm, height=6.5cm,
    ymin=-1000, ymax=1600, ytick={-1000,-500,...,1500},
    tick label style={font=\footnotesize}, label style={font=\small},
    legend style={font=\footnotesize,at={(1.03,1)},anchor=north west,draw=none},
    legend cell align=left]
  \addplot[Ink,line width=1.1pt,domain=0:0.02,samples=241]
    {500 - 1000*cos(36000*x - 60)};
  \addlegendentry{$p(t)$}
  \addplot[CorP,line width=1pt,dashed,domain=0:0.02,samples=241]
    {500*(1 - cos(36000*x))};
  \addlegendentry{$P\,(1 - \cos 2\omega t)$}
  \addplot[CorS,line width=1pt,dash dot,domain=0:0.02,samples=241]
    {-866*sin(36000*x)};
  \addlegendentry{$-Q\sen 2\omega t$}
\end{axis}
\end{tikzpicture}}

% ---------------------------------------------------------------------
\hypersetup{pdftitle={Resolução detalhada --- Questões 51, 24 e 45},
            pdfauthor={Luis Felipe Carneiro de Souza}}

\begin{document}

\thispagestyle{plain}
\begin{center}
  \includegraphics[height=1.5cm]{../../TemplateBeamerUFC/libs/grei_full_logo.pdf}\par
  \vspace{1em}
  {\LARGE\bfseries\color{blue_theme} Resolução detalhada de exercícios\par}
  \vspace{0.4em}
  {\large Sistemas Elétricos de Potência --- Questões 51, 24 e 45\par}
  \vspace{0.8em}
  Luis Felipe Carneiro de Souza\par
  {\small GREI -- Grupo de Redes Elétricas Inteligentes\\
   Universidade Federal do Ceará\par}
\end{center}

\vspace{0.5em}
\begin{observacao}[title=Sobre este documento]
  Este texto acompanha a apresentação de slides sobre as mesmas questões. Nos
  slides, boa parte da explicação é dada oralmente; aqui ela vem por extenso:
  cada conteúdo usado é revisado com a dedução das fórmulas, e cada conta da
  resolução é feita passo a passo, com as conferências. A organização é a
  mesma da apresentação. Para cada questão: o enunciado, os conteúdos
  envolvidos, a revisão desses conteúdos e a resolução.
\end{observacao}

\tableofcontents

\subsection*{Notação e convenções}
\begin{itemize}
  \item Sistema trifásico equilibrado, representado por fase no diagrama
        unifilar. Tensões são \emph{de linha} e potências são
        \emph{trifásicas}, salvo aviso.
  \item Grandezas em pu (por unidade) não têm unidade. O separador decimal é a
        vírgula: $0{,}5$~pu.
  \item Fasores: $V = |V|\angle\theta = |V|\cos\theta + j\,|V|\sin\theta$,
        com $j = \sqrt{-1}$. O asterisco indica o conjugado:
        $V^{*} = |V|\angle{-\theta}$.
  \item $\operatorname{Re}\{\cdot\}$ e $\operatorname{Im}\{\cdot\}$ são as
        partes real e imaginária. Ângulos em graus.
  \item Linhas e transformadores sem perdas: representados só pela reatância
        série.
\end{itemize}

% =====================================================================
\section{Questão 51 --- Perda de uma linha de transmissão (ENADE 2008)}
% =====================================================================

\subsection{Enunciado}

\begin{center}
  \unifilarQcinquentaum
\end{center}

\begin{center}
\begin{minipage}{0.75\textwidth}
  Dados dos componentes:
  \begin{itemize}[nosep]
    \item Gerador $G_1$: 100 MVA, 20 kV, $X = 10\%$
    \item Gerador $G_2$: 100 MVA, 25 kV, $X = 10\%$
    \item Transformadores $T_1$ e $T_2$: 50 MVA, 20 kV/250 kV, $X = 10\%$
    \item Transformadores $T_3$ e $T_4$: 50 MVA, 25 kV/250 kV, $X = 10\%$
  \end{itemize}
\end{minipage}
\end{center}

A figura acima apresenta um sistema de potência sem perdas. Suponha que,
inicialmente, a potência ativa transmitida do gerador $G_1$ para o $G_2$ seja
igual a 1,0~p.u., e que os dois estejam operando em suas tensões nominais. Em
um dado instante, uma das linhas de transmissão é perdida. Supondo que as
tensões dos geradores e a abertura angular permaneçam inalteradas, qual será a
nova potência ativa transmitida, em p.u., nessa condição de regime?

\subsection{Conteúdos envolvidos}

A questão pede a potência transmitida depois de uma mudança na rede: uma das
duas linhas sai de operação, enquanto as tensões dos geradores e a abertura
angular entre eles ficam como estavam. Ela envolve conceitos de
\textbf{estabilidade de SEP} e de \textbf{sistema por unidade}. Três
conteúdos são necessários:

\begin{enumerate}
  \item \passo{Sistema por unidade (pu).} Os dados de cada equipamento estão
        em porcentagem na base do próprio equipamento: potência nominal e
        tensão nominal. Além disso, o sistema tem três níveis de tensão (20,
        250 e 25~kV). Para montar um único circuito, todas as reatâncias
        precisam estar na mesma base.
  \item \passo{Associação de reatâncias.} Os transformadores e as linhas
        estão em paralelo, e cada par está em série com os demais e com os
        geradores. Tudo isso se reduz a uma reatância equivalente $X_{eq}$
        entre os dois geradores.
  \item \passo{Equação potência-ângulo.} É a relação, estudada em
        estabilidade de SEP, entre a potência ativa transmitida entre duas
        fontes, as tensões delas, a abertura angular $\delta$ e a reatância
        entre elas. É ela que diz como a potência muda quando $X_{eq}$ muda.
\end{enumerate}

\subsection{Revisão do conteúdo}

\subsubsection{Sistema por unidade}

Num sistema de potência convivem tensões de poucos kV a centenas de kV e
potências de kVA a GVA. Expressar cada grandeza como uma fração de um valor
de referência, a \emph{base}, deixa os números na mesma ordem de grandeza e
simplifica os cálculos \cite{stevenson1986elementos,bichels2018sep}:
\begin{equation}
  \text{valor em pu} = \frac{\text{valor real}}{\text{valor base}}.
  \label{eq:pu}
\end{equation}
O valor em porcentagem é o valor em pu multiplicado por 100. Assim, uma
reatância de $X = 10\,\%$ vale $0{,}10$~pu \emph{na base em que foi dada}.

\paragraph{Escolha das bases.}
As quatro grandezas de interesse (potência $S$, tensão $V$, corrente $I$ e
impedância $Z$) estão ligadas por duas relações: $S = \sqrt{3}\,V I$ e,
por fase, $Z = V_{\text{fase}}/I$. Por isso só duas delas podem ser
escolhidas livremente. Na prática escolhem-se:
\begin{itemize}
  \item a potência base $S_b$ (trifásica), \emph{uma só para todo o sistema};
  \item a tensão base $V_b$ (de linha), \emph{uma para cada zona} (ver
        adiante).
\end{itemize}
As outras duas bases saem dessas. A corrente base vem de
$S_b = \sqrt{3}\,V_b I_b$:
\begin{equation}
  I_b = \frac{S_b}{\sqrt{3}\,V_b}.
  \label{eq:Ib}
\end{equation}
A impedância base é a tensão base de fase dividida pela corrente base:
\begin{equation}
  Z_b = \frac{V_b/\sqrt{3}}{I_b}
      = \frac{V_b}{\sqrt{3}}\cdot\frac{\sqrt{3}\,V_b}{S_b}
      = \frac{V_b^{2}}{S_b}.
  \label{eq:Zb}
\end{equation}
Com $V_b$ em kV e $S_b$ em MVA, a Equação~\eqref{eq:Zb} já dá $Z_b$ em ohms,
pois $(\text{kV})^2/\text{MVA} = 10^{6}\,\text{V}^2/10^{6}\,\text{VA} =
\Omega$.

\paragraph{Zonas e transformadores.}
Considere um transformador ideal de tensões nominais $V_{1n}/V_{2n}$ e uma
impedância $Z_2$ (em ohms) no lado 2. Refletida para o lado 1, ela vale
$Z_2' = Z_2\,(V_{1n}/V_{2n})^{2}$. Escolhendo as tensões base dos dois lados
na mesma relação do transformador, $V_{b1}/V_{b2} = V_{1n}/V_{2n}$, o valor
em pu dessa impedância é o mesmo visto de qualquer lado:
\begin{equation}
  \frac{Z_2'}{Z_{b1}}
  = \frac{Z_2\,(V_{1n}/V_{2n})^{2}}{V_{b1}^{2}/S_b}
  = \frac{Z_2\,(V_{b1}/V_{b2})^{2}}{V_{b1}^{2}/S_b}
  = \frac{Z_2}{V_{b2}^{2}/S_b}
  = \frac{Z_2}{Z_{b2}}.
  \label{eq:trafo-pu}
\end{equation}
Ou seja, em pu o transformador ideal passa a ter relação 1:1 e
\emph{desaparece} do circuito; sobra apenas a sua reatância de dispersão.
Por isso a tensão base muda ao atravessar cada transformador, na relação de
transformação dele. Cada região delimitada por transformadores tem uma única
tensão base e é chamada de \emph{zona}.

\paragraph{Mudança de base.}
Os fabricantes informam a reatância de cada equipamento em pu (ou \%) na
base do próprio equipamento: a sua potência e a sua tensão nominais. O valor
em ohms, porém, é um só, qualquer que seja a base:
\[
  Z_{\Omega} = Z_{\text{pu,velho}}\,Z_{b,\text{velho}}
             = Z_{\text{pu,novo}}\,Z_{b,\text{novo}} .
\]
Isolando $Z_{\text{pu,novo}}$ e usando a Equação~\eqref{eq:Zb}:
\[
  Z_{\text{pu,novo}}
  = Z_{\text{pu,velho}}\,\frac{Z_{b,\text{velho}}}{Z_{b,\text{novo}}}
  = Z_{\text{pu,velho}}\,
    \frac{V_{b,\text{velho}}^{2}/S_{b,\text{velho}}}
         {V_{b,\text{novo}}^{2}/S_{b,\text{novo}}},
\]
que, reorganizada, dá a fórmula de mudança de base:
\begin{conceito}[title=Mudança de base]
  \begin{equation}
    Z_{\text{pu,novo}} = Z_{\text{pu,velho}}
    \left(\frac{S_{b,\text{novo}}}{S_{b,\text{velho}}}\right)
    \left(\frac{V_{b,\text{velho}}}{V_{b,\text{novo}}}\right)^{2}
    \label{eq:mudanca-base}
  \end{equation}
\end{conceito}
Quando a tensão nominal do equipamento coincide com a tensão base da zona
onde ele está ($V_{b,\text{velho}} = V_{b,\text{novo}}$), o último fator
vale 1 e só a relação de potências altera o valor.

\subsubsection{Associação de reatâncias}

Reatâncias se associam como resistências. Em série, somam-se:
\begin{equation}
  X_s = X_1 + X_2 + \dots + X_n .
  \label{eq:serie}
\end{equation}
Em paralelo, somam-se os inversos:
\begin{equation}
  \frac{1}{X_p} = \frac{1}{X_1} + \frac{1}{X_2} + \dots + \frac{1}{X_n}
  \quad\Longrightarrow\quad
  X_p = \frac{X_1 X_2}{X_1 + X_2}\ \ \text{(duas reatâncias)}.
  \label{eq:paralelo}
\end{equation}
Duas reatâncias \emph{iguais} em paralelo valem metade de uma delas:
$X_p = X\cdot X/(2X) = X/2$. Em pu, as mesmas regras valem, desde que todas
as reatâncias estejam na mesma base.

\subsubsection{Equação potência-ângulo}

Em estudos de estabilidade, cada máquina síncrona é representada por uma
tensão interna $E$ atrás da sua reatância. Num sistema sem perdas, tudo o
que há entre as tensões internas de duas máquinas (reatâncias das máquinas,
dos transformadores e das linhas) se reduz a uma única reatância $X$
\cite{stevenson1986elementos}. Tomando a tensão da fonte 2 como referência
angular, $E_2\angle 0^\circ$, a da fonte 1 fica adiantada de $\delta$:
$E_1\angle\delta$ (Figura~\ref{fig:duasfontes}).

\begin{figure}[H]
  \centering
  \duasfontes
  \caption{Duas fontes ligadas por uma reatância $X$; $\delta$ é a abertura
    angular entre elas.}
  \label{fig:duasfontes}
\end{figure}

A corrente que vai da fonte 1 para a fonte 2 sai da lei de Ohm na
reatância:
\begin{equation}
  I = \frac{E_1\angle\delta - E_2\angle 0^\circ}{jX}.
  \label{eq:I-duasfontes}
\end{equation}
A potência complexa entregue pela fonte 1 é $S_1 = E_1\angle\delta\cdot
I^{*}$. Conjugando a Equação~\eqref{eq:I-duasfontes} (o conjugado de $jX$
é $-jX$):
\[
  I^{*} = \frac{E_1\angle{-\delta} - E_2}{-jX}.
\]
Multiplicando por $E_1\angle\delta$ e usando $\angle\delta\cdot
\angle{-\delta} = 1$ e $1/(-j) = j$:
\begin{align*}
  S_1 &= \frac{E_1\angle\delta\,\bigl(E_1\angle{-\delta} - E_2\bigr)}{-jX}
       = \frac{E_1^{2} - E_1E_2\angle\delta}{-jX}
       = \frac{j}{X}\,\bigl(E_1^{2} - E_1E_2\angle\delta\bigr).
\end{align*}
Escrevendo $E_1E_2\angle\delta = E_1E_2\cos\delta + j\,E_1E_2\sin\delta$ e
usando $j\cdot j = -1$:
\begin{align*}
  S_1 &= \frac{j}{X}\bigl(E_1^{2} - E_1E_2\cos\delta\bigr)
        - \frac{j\cdot j}{X}\,E_1E_2\sin\delta \\
      &= {\underbrace{\frac{E_1E_2}{X}\sin\delta}_{P}}
        + j\,{\underbrace{\frac{E_1^{2} - E_1E_2\cos\delta}{X}}_{Q_1}}.
\end{align*}
A parte real é a potência ativa transmitida:
\begin{conceito}[title=Equação potência-ângulo]
  \begin{equation}
    P = \frac{E_1\,E_2}{X}\,\sin\delta
    \label{eq:potencia-angulo}
  \end{equation}
\end{conceito}

Repetindo a conta para a potência que chega à fonte 2, $S_2 = E_2\,I^{*}$,
a parte real é a mesma: como não há resistência, toda a potência ativa que
sai de uma fonte chega à outra. Três consequências da
Equação~\eqref{eq:potencia-angulo}:

\begin{itemize}
  \item \passo{Potência máxima.} $\sin\delta$ é máximo em $\delta =
        90^\circ$, então $P_{\max} = E_1E_2/X$. Acima de $90^\circ$, aumentar
        a abertura angular \emph{diminui} a potência, e as máquinas perdem o
        sincronismo. $P_{\max}$ é o limite de estabilidade em regime
        permanente (Figura~\ref{fig:curva-Pdelta}).
  \item \passo{Sentido do fluxo.} $P > 0$ para $0 < \delta < 180^\circ$: a
        potência ativa flui da fonte adiantada para a atrasada.
  \item \passo{Papel da reatância.} Com $E_1$, $E_2$ e $\delta$ fixos, $P$ é
        inversamente proporcional a $X$: aumentar a reatância entre as
        fontes reduz a potência transmitida na mesma proporção.
\end{itemize}

\begin{figure}[H]
  \centering
  \curvaPdelta
  \caption{Curva potência-ângulo $P = P_{\max}\sin\delta$.}
  \label{fig:curva-Pdelta}
\end{figure}

\subsection{Resolução}

\subsubsection{Leitura do enunciado e hipóteses}

\begin{itemize}
  \item \passo{\enquote{Sistema sem perdas}:} todos os elementos são
        reatâncias puras. A potência que sai de $G_1$ chega inteira a $G_2$.
  \item \passo{\enquote{Operando em suas tensões nominais}:} $E_1 = E_2 =
        1{,}0$~pu, pois, como se verá no Passo~1, as tensões base escolhidas
        coincidem com as nominais dos geradores.
  \item \passo{\enquote{Tensões dos geradores e abertura angular
        inalteradas}:} $E_1$, $E_2$ e $\delta$ são os mesmos antes e depois
        da perda da linha. Adota-se que são as tensões \emph{internas} das
        máquinas, atrás das reatâncias $X_{G}$, que é a forma usual da
        equação potência-ângulo. Por isso $X_{G}$ faz parte da reatância entre
        as fontes. A outra leitura possível é discutida na
        Seção~\ref{sec:q51-conferencias}.
  \item \passo{\enquote{Uma das linhas é perdida}:} as duas linhas são
        iguais, então tanto faz qual delas sai; a outra continua em operação.
\end{itemize}

\subsubsection{Estratégia}

Aplicando a Equação~\eqref{eq:potencia-angulo} antes e depois da perda da
linha:
\[
  P = \frac{E_1E_2}{X_{eq}}\sin\delta
  \qquad\text{e}\qquad
  P' = \frac{E_1E_2}{X_{eq}'}\sin\delta .
\]
Dividindo uma pela outra, $E_1$, $E_2$ e $\delta$ (que não mudam) se
cancelam:
\begin{equation}
  \frac{P'}{P} = \frac{X_{eq}}{X_{eq}'}
  \quad\Longrightarrow\quad
  P' = P\,\frac{X_{eq}}{X_{eq}'} .
  \label{eq:razao}
\end{equation}
Não é preciso conhecer os valores de $E_1$, $E_2$ e $\delta$, apenas as
reatâncias equivalentes antes ($X_{eq}$) e depois ($X_{eq}'$), as duas na
mesma base. O roteiro é:
\begin{enumerate}[label=\arabic*.]
  \item identificar as zonas e escolher as bases;
  \item passar todas as reatâncias para pu na base comum;
  \item montar o circuito equivalente e calcular $X_{eq}$;
  \item repetir sem uma das linhas e calcular $X_{eq}'$;
  \item aplicar a Equação~\eqref{eq:razao}.
\end{enumerate}

\subsubsection{Passo 1: zonas e bases}

Os transformadores dividem o sistema em três zonas
(Figura~\ref{fig:zonas}). A fronteira entre duas zonas passa \emph{por
dentro} de cada transformador: cada enrolamento fica na zona da sua tensão
nominal.

\begin{figure}[H]
  \centering
  \zonasQcinquentaum
  \caption{Zonas de tensão base da Questão 51. As fronteiras (tracejadas)
    passam pelo meio dos transformadores.}
  \label{fig:zonas}
\end{figure}

\passo{Potência base.} Escolhe-se $S_b = 100$~MVA, que é a potência nominal
dos dois geradores. Ela vale para todo o sistema.

\passo{Tensões base.} Escolhe-se a tensão base da zona~1 igual à tensão
nominal de $G_1$, $V_{b1} = 20$~kV. As outras saem das relações de
transformação:
\begin{align*}
  \text{zona 2 (atravessando $T_1$, $T_2$: 20/250 kV):}&\quad
  V_{b2} = 20 \times \frac{250}{20} = 250~\text{kV},\\
  \text{zona 3 (atravessando $T_3$, $T_4$: 250/25 kV):}&\quad
  V_{b3} = 250 \times \frac{25}{250} = 25~\text{kV}.
\end{align*}
A tensão base da zona~3 coincide com a tensão nominal de $G_2$ (25~kV). Isso
confirma que as relações dos transformadores são compatíveis entre si.

\passo{Impedâncias e correntes base.} Pelas Equações~\eqref{eq:Zb}
e~\eqref{eq:Ib}, com $S_b = 100$~MVA:

\begin{table}[H]
  \centering
  \begin{tabular}{@{}lccc@{}}
    \toprule
    Zona & $V_b$ (kV) & $Z_b = V_b^2/S_b$ ($\Omega$) & $I_b = S_b/(\sqrt{3}\,V_b)$ (kA) \\
    \midrule
    1 ($G_1$)            & 20  & $20^2/100 = 4$       & $100/(\sqrt{3}\cdot 20) = 2{,}887$ \\
    2 (linhas)           & 250 & $250^2/100 = 625$    & $100/(\sqrt{3}\cdot 250) = 0{,}231$ \\
    3 ($G_2$)            & 25  & $25^2/100 = 6{,}25$  & $100/(\sqrt{3}\cdot 25) = 2{,}309$ \\
    \bottomrule
  \end{tabular}
  \caption{Bases de cada zona ($S_b = 100$~MVA).}
  \label{tab:bases}
\end{table}

Para esta questão só $Z_{b2} = 625~\Omega$ é necessário, porque é a única
impedância dada em ohms (a das linhas). As correntes base ficam como
referência.

\subsubsection{Passo 2: reatâncias em pu na base comum}

Aplica-se a Equação~\eqref{eq:mudanca-base} a cada equipamento. A
\enquote{base velha} é a do próprio equipamento (seus dados de placa); a
\enquote{base nova} é a do sistema ($S_b = 100$~MVA e a $V_b$ da zona
onde ele está).

\passo{Geradores.} $G_1$: $X = 0{,}10$~pu na base 100~MVA, 20~kV, na
zona~1 ($V_b = 20$~kV):
\[
  X_{G1} = 0{,}10 \times \frac{100}{100} \times
           \left(\frac{20}{20}\right)^{2} = 0{,}10~\text{pu}.
\]
$G_2$: $X = 0{,}10$~pu na base 100~MVA, 25~kV, na zona~3 ($V_b = 25$~kV):
\[
  X_{G2} = 0{,}10 \times \frac{100}{100} \times
           \left(\frac{25}{25}\right)^{2} = 0{,}10~\text{pu}.
\]
Os geradores já estavam na base do sistema, e os valores não mudam.

\passo{Transformadores.} $T_1$ e $T_2$: $X = 0{,}10$~pu na base 50~MVA,
20/250~kV. Pelo lado de 20~kV (zona~1):
\[
  X_{T1} = X_{T2} = 0{,}10 \times \frac{100}{50} \times
           \left(\frac{20}{20}\right)^{2} = 0{,}20~\text{pu}.
\]
Pelo lado de 250~kV (zona~2), o resultado é o mesmo, como garante a
Equação~\eqref{eq:trafo-pu}:
\[
  0{,}10 \times \frac{100}{50} \times
  \left(\frac{250}{250}\right)^{2} = 0{,}20~\text{pu}.
\]
$T_3$ e $T_4$: $X = 0{,}10$~pu na base 50~MVA, 25/250~kV. Pelo lado de
25~kV (zona~3):
\[
  X_{T3} = X_{T4} = 0{,}10 \times \frac{100}{50} \times
           \left(\frac{25}{25}\right)^{2} = 0{,}20~\text{pu}.
\]
Aqui só a potência muda de base: o transformador de 50~MVA visto numa base
de 100~MVA tem reatância em pu duas vezes maior.

\passo{Linhas.} A reatância de cada linha é dada em ohms, $X_L =
125~\Omega$, na zona~2. Pela Equação~\eqref{eq:pu}:
\[
  X_{L1} = X_{L2} = \frac{125~\Omega}{Z_{b2}}
                  = \frac{125}{625} = 0{,}20~\text{pu}.
\]

\passo{Conferência em ohms.} A reatância de $T_1$ vista do lado de 250~kV,
em ohms, é $0{,}10 \times Z_{b,T}$, com $Z_{b,T} = 250^2/50 = 1250~\Omega$
(base do próprio transformador), ou seja, $125~\Omega$. Em pu do sistema:
$125/625 = 0{,}20$, o mesmo valor obtido acima. Coincidentemente, cada
transformador tem, em ohms do lado de 250~kV, a mesma reatância de uma
linha.

\begin{table}[H]
  \centering
  \begin{tabular}{@{}llc@{}}
    \toprule
    Elemento & Cálculo & $X$ (pu, base 100 MVA) \\
    \midrule
    $G_1$, $G_2$       & $0{,}10 \times (100/100)$     & $0{,}10$ \\
    $T_1$ a $T_4$      & $0{,}10 \times (100/50)$      & $0{,}20$ \\
    $L_1$, $L_2$       & $125/625$                     & $0{,}20$ \\
    \bottomrule
  \end{tabular}
  \caption{Reatâncias na base comum.}
  \label{tab:reatancias}
\end{table}

\subsubsection{Passo 3: reatância equivalente antes da perda da linha}

Com todas as reatâncias em pu, os transformadores ideais desaparecem e o
sistema vira o circuito da Figura~\ref{fig:equiv-antes}.

\begin{figure}[H]
  \centering
  \equivQcinquentaum{0}
  \caption{Circuito equivalente em pu, com as duas linhas em operação.}
  \label{fig:equiv-antes}
\end{figure}

Cada par em paralelo tem duas reatâncias iguais de $0{,}20$~pu. Pela
Equação~\eqref{eq:paralelo}:
\[
  X_{T1}\parallel X_{T2}
  = \frac{0{,}20 \times 0{,}20}{0{,}20 + 0{,}20}
  = \frac{0{,}04}{0{,}40} = 0{,}10~\text{pu}.
\]
Da mesma forma, $X_{L1}\parallel X_{L2} = 0{,}10$~pu e $X_{T3}\parallel
X_{T4} = 0{,}10$~pu. Os três pares estão em série entre si e com as
reatâncias dos geradores. Pela Equação~\eqref{eq:serie}:
\begin{equation}
  X_{eq} = {\underbrace{0{,}10}_{G_1}}
         + {\underbrace{0{,}10}_{T_1\parallel T_2}}
         + {\underbrace{0{,}10}_{L_1\parallel L_2}}
         + {\underbrace{0{,}10}_{T_3\parallel T_4}}
         + {\underbrace{0{,}10}_{G_2}}
         = 0{,}50~\text{pu}.
  \label{eq:Xeq}
\end{equation}

\subsubsection{Passo 4: reatância equivalente depois da perda da linha}

Com uma das linhas aberta (Figura~\ref{fig:equiv-depois}), o bloco das
linhas deixa de ser um paralelo: resta uma única linha de $0{,}20$~pu. Os
demais blocos não mudam.

\begin{figure}[H]
  \centering
  \equivQcinquentaum{1}
  \caption{Circuito equivalente em pu, com a linha 2 fora de operação.}
  \label{fig:equiv-depois}
\end{figure}

\begin{equation}
  X_{eq}' = 0{,}10 + 0{,}10 + {\underbrace{0{,}20}_{L_1}} + 0{,}10 + 0{,}10
          = 0{,}60~\text{pu}.
  \label{eq:Xeq-linha}
\end{equation}

\subsubsection{Passo 5: nova potência transmitida}

Substituindo $P = 1{,}0$~pu e as Equações~\eqref{eq:Xeq}
e~\eqref{eq:Xeq-linha} na Equação~\eqref{eq:razao}:
\[
  P' = P\,\frac{X_{eq}}{X_{eq}'}
     = 1{,}0 \times \frac{0{,}50}{0{,}60}
     = \frac{5}{6}
     \approx 0{,}833~\text{pu}.
\]
Em unidades físicas, com $S_b = 100$~MVA, a potência transmitida cai de
$1{,}0 \times 100 = 100$~MW para $0{,}833 \times 100 \approx 83{,}3$~MW.

\begin{resposta}
  \centering
  $P' = \dfrac{5}{6} \approx 0{,}833$~p.u.\quad (cerca de 83,3~MW na base de
  100~MVA)
\end{resposta}

\subsubsection{Conferências e interpretação}
\label{sec:q51-conferencias}

\passo{Conferência com os valores absolutos.} Com $E_1 = E_2 = 1{,}0$~pu, a
abertura angular inicial sai da Equação~\eqref{eq:potencia-angulo}:
\[
  \sin\delta = \frac{P\,X_{eq}}{E_1E_2} = \frac{1{,}0 \times 0{,}50}{1{,}0
  \times 1{,}0} = 0{,}5
  \quad\Longrightarrow\quad
  \delta = 30^\circ .
\]
Mantendo $\delta = 30^\circ$ com a nova reatância:
\[
  P' = \frac{1{,}0 \times 1{,}0}{0{,}60}\,\sin 30^\circ
     = \frac{0{,}5}{0{,}6} \approx 0{,}833~\text{pu},
\]
o mesmo resultado do Passo~5.

\passo{Leitura nas curvas potência-ângulo.} Antes, $P_{\max} =
1/0{,}50 = 2{,}0$~pu; depois, $P_{\max}' = 1/0{,}60 \approx 1{,}667$~pu. A
perda da linha \enquote{achata} a curva (Figura~\ref{fig:curvas-q51}). Com
$\delta$ fixo em $30^\circ$, o ponto de operação desce da curva de antes
($1{,}0$~pu) para a de depois ($0{,}833$~pu).

\begin{figure}[H]
  \centering
  \curvasQcinquentaum
  \caption{Curvas potência-ângulo antes e depois da perda da linha, com o
    ponto de operação em $\delta = 30^\circ$.}
  \label{fig:curvas-q51}
\end{figure}

\passo{Por que a potência não cai pela metade?} Perde-se metade das linhas,
mas as linhas são só uma parte da reatância total. Antes, o paralelo das
linhas contribuía com $0{,}10$ dos $0{,}50$~pu; depois, a linha restante
contribui com $0{,}20$. A reatância total sobe de $0{,}50$ para $0{,}60$
(20\,\%), e a potência cai na proporção inversa, para $5/6$ do valor inicial
(uma queda de 16,7\,\%).

\begin{observacao}[title=Observação: o que aconteceria na prática]
  O enunciado fixa a abertura angular. Num sistema real, logo após a perda
  da linha a potência mecânica das turbinas continua a mesma, e o rotor de
  $G_1$ acelera até que $\delta$ aumente o bastante para a rede voltar a
  transmitir $1{,}0$~pu. O novo regime teria $\sin\delta' = 1{,}0 \times
  0{,}60 = 0{,}6$, ou seja, $\delta' \approx 36{,}9^\circ$, ainda abaixo de
  $90^\circ$ e portanto estável. A questão pede outra coisa: com $\delta$
  mantido, quanto vale a potência.
\end{observacao}

\passo{Outra leitura do enunciado.} Se as tensões mantidas fossem as dos
\emph{terminais} dos geradores (barras de 20 e 25~kV), as reatâncias $X_{G}$
ficariam fora do caminho entre as duas tensões:
\[
  X_{eq} = 0{,}10 + 0{,}10 + 0{,}10 = 0{,}30~\text{pu},
  \qquad
  X_{eq}' = 0{,}10 + 0{,}20 + 0{,}10 = 0{,}40~\text{pu},
\]
\[
  P' = 1{,}0 \times \frac{0{,}30}{0{,}40} = 0{,}75~\text{pu}.
\]
Este texto adota a leitura com as tensões internas ($0{,}833$~pu), que é a
usual quando se fala em abertura angular entre geradores.

% =====================================================================
\newpage
\section{Questão 24 --- Fluxo de potência em três barras}
% =====================================================================

\subsection{Enunciado}

Considere um sistema elétrico de potência representado pelo diagrama
unifilar a seguir. A potência injetada na barra 3 e as impedâncias das
linhas que conectam a barra 1 à barra 2 bem como a barra 2 à barra 3 são
mensuradas por unidade (pu). Tanto a tensão na barra 1 quanto a tensão na
barra 3 são iguais a 1,0~pu.

\begin{center}
  {\small\bfseries Diagrama unifilar do sistema}\\[0.5em]
  \unifilarQvintequatro{0}
\end{center}

Com base nessas informações, é correto dizer que a tensão na barra 2 e o
fluxo de potência saindo da barra 1 para a barra 2 são, respectivamente,
\begin{enumerate}[label=(\Alph*)]
  \item $1{,}0$ pu e $1{,}0$ pu.
  \item $1{,}0-j0{,}1$ pu e $1{,}0$ pu.
  \item $1{,}0+j0{,}1$ pu e $0{,}5$ pu.
  \item $1{,}0-j0{,}1$ pu e $0{,}5$ pu.
  \item $1{,}0+j0{,}1$ pu e $1{,}0$ pu.
\end{enumerate}

\subsection{Conteúdos envolvidos}

A questão é um pequeno problema de \textbf{fluxo de potência}: dadas
algumas tensões e potências, encontrar as que faltam. Ela usa:

\begin{enumerate}
  \item \passo{Sistema por unidade.} Tensões, impedâncias e potência já vêm
        em pu, então as contas são feitas diretamente com esses números (a
        revisão está na Questão~51).
  \item \passo{Fasores e potência complexa.} É a relação entre a potência,
        a tensão e a corrente numa barra, $S = V I^{*}$, que permite passar
        da potência injetada para a corrente.
  \item \passo{Fluxo numa linha.} A queda de tensão na reatância de cada
        linha (lei de Ohm com fasores) e a potência que sai de uma barra para
        a outra.
\end{enumerate}

\subsection{Revisão do conteúdo}

\subsubsection{Fasores e potência complexa}

Uma tensão senoidal $v(t) = \sqrt{2}\,|V|\cos(\omega t + \theta_V)$ é
representada pelo fasor $V = |V|\angle\theta_V$, e uma corrente, da mesma
forma, por $I = |I|\angle\theta_I$. A \emph{potência complexa} numa barra é
\begin{equation}
  S = V\,I^{*} = |V|\,|I|\angle(\theta_V - \theta_I) = P + jQ,
  \label{eq:S}
\end{equation}
com
\[
  P = |V|\,|I|\cos\varphi, \qquad Q = |V|\,|I|\sin\varphi, \qquad
  \varphi = \theta_V - \theta_I .
\]
O conjugado em $I^{*}$ aparece justamente para que o ângulo de $S$ seja a
\emph{diferença} entre os ângulos da tensão e da corrente, isto é, o ângulo
do fator de potência \cite{stevenson1986elementos}. Conhecidas a potência e
a tensão, a corrente sai invertendo a Equação~\eqref{eq:S}:
\begin{conceito}[title=Corrente a partir da potência]
  \begin{equation}
    S = V\,I^{*}
    \quad\Longrightarrow\quad
    I^{*} = \frac{S}{V}
    \quad\Longrightarrow\quad
    I = \left(\frac{S}{V}\right)^{\!*}
    \label{eq:I-de-S}
  \end{equation}
\end{conceito}
Em pu, a mesma relação vale para o sistema trifásico: o fator $\sqrt{3}$ de
$S = \sqrt{3}\,V I$ é absorvido pelas bases (Equação~\eqref{eq:Ib}). Por
convenção, $S$ é a potência \emph{injetada} na rede pela barra: positiva
para um gerador e negativa para uma carga.

\subsubsection{Fluxo de potência numa linha}

Uma linha curta e sem perdas entre as barras $k$ e $m$ é representada só
pela sua reatância série $x$ (Figura~\ref{fig:linhaKM})
\cite{stevenson1986elementos,bichels2018sep,borges2005analise}.

\begin{figure}[H]
  \centering
  \linhaKM
  \caption{Linha entre as barras $k$ e $m$, só com reatância.}
  \label{fig:linhaKM}
\end{figure}

A lei de Ohm na reatância dá a queda de tensão e, daí, a corrente:
\begin{equation}
  V_k - V_m = jx\,I_{km}
  \quad\Longrightarrow\quad
  I_{km} = \frac{V_k - V_m}{jx}.
  \label{eq:Ikm}
\end{equation}
A potência que sai da barra $k$ em direção a $m$ é
\begin{equation}
  S_{km} = V_k\,I_{km}^{*}.
  \label{eq:Skm}
\end{equation}
Substituindo a Equação~\eqref{eq:Ikm}, com $V_k = |V_k|\angle\theta_k$,
$V_m = |V_m|\angle\theta_m$ e $\theta_{km} = \theta_k - \theta_m$:
\[
  S_{km} = V_k\,\frac{V_k^{*} - V_m^{*}}{-jx}
         = \frac{j}{x}\bigl(|V_k|^{2} - V_k V_m^{*}\bigr),
  \qquad
  V_k V_m^{*} = |V_k|\,|V_m|\bigl(\cos\theta_{km} + j\sin\theta_{km}\bigr).
\]
Separando parte real e parte imaginária, como na dedução da equação
potência-ângulo:
\begin{conceito}[title=Fluxo numa linha sem perdas]
  \begin{align}
    P_{km} &= \frac{|V_k|\,|V_m|}{x}\,\sin\theta_{km}, \label{eq:Pkm}\\
    Q_{km} &= \frac{|V_k|^{2} - |V_k|\,|V_m|\cos\theta_{km}}{x}. \label{eq:Qkm}
  \end{align}
\end{conceito}

Propriedades úteis:
\begin{itemize}
  \item A Equação~\eqref{eq:Pkm} é a equação potência-ângulo da
        Questão~51, com as tensões das barras no lugar das tensões internas.
        A potência ativa vai da barra de \emph{maior ângulo} para a de menor.
  \item Sem resistência não há perda ativa: $P_{mk} = -P_{km}$.
  \item A linha consome potência reativa. Somando $Q_{km}$ e $Q_{mk}$:
        \[
          Q_{km} + Q_{mk}
          = \frac{|V_k|^{2} + |V_m|^{2} - 2\,|V_k|\,|V_m|\cos\theta_{km}}{x}
          = \frac{|V_k - V_m|^{2}}{x}
          = |I_{km}|^{2}\,x ,
        \]
        em que a última igualdade vem de $|V_k - V_m| = x\,|I_{km}|$
        (Equação~\eqref{eq:Ikm}).
\end{itemize}

\subsection{Resolução}

\subsubsection{Leitura dos dados}

\begin{itemize}
  \item \passo{Tensões.} \enquote{Tanto a tensão na barra 1 quanto a tensão
        na barra 3 são iguais a 1,0~pu} é lido como fasores iguais:
        $V_1 = V_3 = 1{,}0\angle 0^\circ$~pu. A barra 3 pode ser tomada como
        referência angular ($0^\circ$) sem perda de generalidade; que a
        barra 1 também esteja em $0^\circ$ é uma interpretação, compatível
        com as alternativas, que dão $V_2$ com parte real $1{,}0$.
  \item \passo{Potência injetada.} O diagrama traz $S_{G3} = 0{,}5$, um
        número real: a barra 3 injeta $S_{G3} = 0{,}5 + j0$~pu, só potência
        ativa.
  \item \passo{Linhas.} $x = 0{,}2$~pu em cada trecho (1--2 e 2--3).
  \item \passo{Incógnitas.} A tensão $V_2$ e o fluxo $S_{12}$ da barra 1
        para a barra 2. A carga da barra 2 \emph{não} é dada.
\end{itemize}

\begin{observacao}[title={Por que a leitura $V_1 = V_3 = 1\angle 0^\circ$ é necessária}]
  Sem ela, o problema fica indeterminado. Se só os módulos $|V_1| = |V_3| =
  1{,}0$ fossem conhecidos, faltariam o ângulo de $V_1$ e a carga da
  barra~2, que não aparece no enunciado. Haveria infinitas soluções.
\end{observacao}

\subsubsection{Estratégia}

A carga da barra 2 é desconhecida, mas na \emph{barra 3} conhecemos a
tensão \emph{e} a potência. O caminho é:
\begin{enumerate}[label=\arabic*.]
  \item na barra 3, da potência e da tensão sai a corrente $I_{32}$
        (Equação~\eqref{eq:I-de-S});
  \item a queda na linha 3--2 dá $V_2$ (Equação~\eqref{eq:Ikm});
  \item com $V_1$ e $V_2$, a corrente $I_{12}$ e o fluxo $S_{12}$
        (Equações~\eqref{eq:Ikm} e~\eqref{eq:Skm}).
\end{enumerate}

\subsubsection{Passo 1: corrente injetada na barra 3}

Da barra 3 sai uma única linha, a 3--2; portanto toda a corrente do
gerador vai para ela: $I_{32} = I_3$. Pela Equação~\eqref{eq:I-de-S}:
\[
  I_{32} = \left(\frac{S_{G3}}{V_3}\right)^{\!*}
         = \left(\frac{0{,}5 + j0}{1{,}0\angle 0^\circ}\right)^{\!*}
         = (0{,}5)^{*}
         = 0{,}5\angle 0^\circ~\text{pu}.
\]
Como a potência é puramente ativa e $V_3$ está em $0^\circ$, a corrente fica
em fase com a tensão.

\subsubsection{Passo 2: tensão na barra 2}

Na linha 3--2, a Equação~\eqref{eq:Ikm} com $k = 3$ e $m = 2$ dá $V_3 - V_2
= jx\,I_{32}$. Isolando $V_2$:
\begin{align*}
  V_2 &= V_3 - jx\,I_{32} \\
      &= 1{,}0 - j0{,}2 \times 0{,}5 \\
      &= 1{,}0 - j0{,}1~\text{pu}.
\end{align*}
Em módulo e ângulo:
\[
  |V_2| = \sqrt{1{,}0^{2} + 0{,}1^{2}} = \sqrt{1{,}01} \approx 1{,}005~\text{pu},
  \qquad
  \theta_2 = \arctan\!\left(\frac{-0{,}1}{1{,}0}\right) \approx -5{,}71^\circ,
\]
ou seja, $V_2 \approx 1{,}005\angle{-5{,}71^\circ}$~pu.

\subsubsection{Passo 3: fluxo da barra 1 para a barra 2}

Na linha 1--2, pela Equação~\eqref{eq:Ikm}:
\[
  I_{12} = \frac{V_1 - V_2}{jx}
         = \frac{1{,}0 - (1{,}0 - j0{,}1)}{j0{,}2}
         = \frac{j0{,}1}{j0{,}2}
         = 0{,}5~\text{pu}.
\]
Pela Equação~\eqref{eq:Skm}:
\[
  S_{12} = V_1\,I_{12}^{*} = 1{,}0\angle 0^\circ \times 0{,}5
         = 0{,}5 + j0~\text{pu}.
\]
Portanto $V_2 = 1{,}0 - j0{,}1$~pu e o fluxo da barra 1 para a barra 2 é
$0{,}5$~pu (só potência ativa).

\begin{resposta}
  \centering
  Alternativa \textbf{(D)}:\quad $V_2 = 1{,}0 - j0{,}1$~pu \quad e \quad
  $S_{12} = 0{,}5$~pu
\end{resposta}

\subsubsection{Conferências}

\passo{Pelas fórmulas em módulo e ângulo.} Com $\theta_{12} = \theta_1 -
\theta_2 = 0 - (-5{,}71^\circ) = 5{,}71^\circ$, as
Equações~\eqref{eq:Pkm} e~\eqref{eq:Qkm} dão:
\[
  P_{12} = \frac{1{,}0 \times 1{,}005}{0{,}2}\,\sin 5{,}71^\circ
         = \frac{0{,}1}{0{,}2} = 0{,}5~\text{pu},
\]
\[
  Q_{12} = \frac{1{,}0^{2} - 1{,}0 \times 1{,}005\cos 5{,}71^\circ}{0{,}2}
         = \frac{1{,}0 - 1{,}0}{0{,}2} = 0 .
\]
As duas contas ficam exatas porque $|V_2|\sin 5{,}71^\circ =
-\operatorname{Im}\{V_2\} = 0{,}1$ e $|V_2|\cos 5{,}71^\circ =
\operatorname{Re}\{V_2\} = 1{,}0$.

\passo{Balanço de potência na barra 2.} As correntes que saem da barra 2
pelas linhas são $I_{21} = -I_{12} = -0{,}5$ e $I_{23} = -I_{32} = -0{,}5$.
Pela Equação~\eqref{eq:Skm}:
\[
  S_{21} = V_2\,I_{21}^{*} = (1{,}0 - j0{,}1)(-0{,}5) = -0{,}5 + j0{,}05,
  \qquad
  S_{23} = -0{,}5 + j0{,}05 .
\]
A injeção total da barra 2 na rede é $S_{21} + S_{23} = -1{,}0 + j0{,}1$.
Logo a carga da barra 2 absorve
\[
  S_c = -(S_{21} + S_{23}) = 1{,}0 - j0{,}1~\text{pu}.
\]
Ela consome $1{,}0$~pu de potência ativa e \emph{fornece} $0{,}1$~pu de
reativo: comporta-se como uma carga capacitiva. O balanço fecha:
\begin{itemize}
  \item \emph{ativa}: $G_1$ e $G_3$ entregam $0{,}5$ cada; a carga consome
        $1{,}0$; as linhas não têm perdas;
  \item \emph{reativa}: cada linha consome $|I|^{2}x = 0{,}5^{2} \times
        0{,}2 = 0{,}05$, ou $0{,}1$ no total. Os geradores não entregam
        reativo ($Q_{12} = Q_{32} = 0$), então esses $0{,}1$ vêm da barra 2.
        Na linha 1--2: $Q_{12} + Q_{21} = 0 + 0{,}05 = 0{,}05 = |I_{12}|^{2}x$.
\end{itemize}

\passo{Eliminação das alternativas.} Com $V_1 = 1\angle 0^\circ$ e $x =
0{,}2$, o fluxo $S_{12}$ depende só de $V_2$:
$S_{12} = V_1\bigl[(V_1 - V_2)/(jx)\bigr]^{*}$. Cada alternativa propõe um
par ($V_2$, $S_{12}$), e basta verificar se ele é coerente:

\begin{table}[H]
  \centering
  \begin{tabular}{@{}cccccc@{}}
    \toprule
    Alt. & $V_2$ proposto & $V_1 - V_2$ & $I_{12} = (V_1 - V_2)/(j0{,}2)$ &
    $S_{12}$ calculado & $S_{12}$ proposto \\
    \midrule
    (A) & $1{,}0$          & $0$         & $0$        & $0$        & $1{,}0$ \ \ \nok \\
    (B) & $1{,}0 - j0{,}1$ & $j0{,}1$    & $0{,}5$    & $0{,}5$    & $1{,}0$ \ \ \nok \\
    (C) & $1{,}0 + j0{,}1$ & $-j0{,}1$   & $-0{,}5$   & $-0{,}5$   & $0{,}5$ \ \ \nok \\
    (D) & $1{,}0 - j0{,}1$ & $j0{,}1$    & $0{,}5$    & $0{,}5$    & $0{,}5$ \ \ \ok \\
    (E) & $1{,}0 + j0{,}1$ & $-j0{,}1$   & $-0{,}5$   & $-0{,}5$   & $1{,}0$ \ \ \nok \\
    \bottomrule
  \end{tabular}
  \caption{Coerência entre $V_2$ e $S_{12}$ em cada alternativa.}
  \label{tab:alternativas}
\end{table}

Só a alternativa (D) é coerente, e isso sem usar o valor de $S_{G3}$. O dado
$S_{G3}$ confirma a escolha: com o $V_2$ de (D), $I_{32} = (V_3 -
V_2)/(jx) = j0{,}1/j0{,}2 = 0{,}5$ e $S_{G3} = V_3\,I_{32}^{*} = 0{,}5$,
exatamente o valor dado. Em (C) e (E), $V_2 = 1 + j0{,}1$ daria $I_{32} =
-0{,}5$ e $S_{G3} = -0{,}5$: a barra 3 estaria \emph{absorvendo} potência,
o oposto do enunciado.

\passo{Diagrama fasorial.} A Figura~\ref{fig:fasorial} mostra as grandezas
em escala. Como $I_{32}$ está em fase com $V_3$, a queda $jx\,I_{32}$ é
perpendicular a $V_3$. Ela altera quase só o \emph{ângulo} de $V_2$
($-5{,}71^\circ$); o módulo varia pouco, de $1{,}0$ para $1{,}005$.

\begin{figure}[H]
  \centering
  \fasorialQvintequatro
  \caption{Diagrama fasorial da Questão 24 (em escala).}
  \label{fig:fasorial}
\end{figure}

\passo{Interpretação física.} O sistema é simétrico: $V_1 = V_3$ e as duas
linhas são iguais. Então cada lado entrega metade da potência ativa da
carga, $0{,}5$~pu, e a barra 2 fica atrasada em ângulo em relação às
outras duas, pois a potência ativa vai do maior para o menor ângulo. A
Figura~\ref{fig:q24-solucao} resume a solução.

\begin{figure}[H]
  \centering
  \unifilarQvintequatro{1}
  \caption{Tensões e fluxos da solução da Questão 24.}
  \label{fig:q24-solucao}
\end{figure}

% =====================================================================
\newpage
\section{Questão 45 --- Potência ativa e reativa a partir da potência
instantânea}
% =====================================================================

\subsection{Enunciado}

Sejam $v(t) = \sqrt{2}\,V\sen(\omega t)$ e $i(t) = \sqrt{2}\,I\sen(\omega t
- \theta)$, respectivamente, a tensão (volts) e a corrente (ampères)
fornecidas por um gerador CA (Corrente Alternada) em regime permanente.

A figura apresenta a curva de potência instantânea fornecida por este
gerador durante o intervalo de tempo $\dfrac{2\pi}{\omega}$.

\begin{center}
  \graficoQquarentacinco[width=12cm,height=7.5cm,
    tick label style={font=\footnotesize},label style={font=\small}]{0}
\end{center}

Denomina-se Potência Ativa $P$ o valor médio da potência instantânea ao
longo do tempo e é dada por:\\
$P = VI\cos\theta$.

Denomina-se Potência Reativa $Q$ o valor calculado pela expressão $Q =
VI\sen\theta$.

Analisando a figura, qual o valor aproximado da Potência Reativa, em var?
\begin{enumerate}[label=(\Alph*)]
  \item 1\,000
  \item 870
  \item 500
  \item 290
  \item Impossível determinar com os dados fornecidos.
\end{enumerate}

\subsection{Conteúdos envolvidos}

A questão dá apenas a curva da potência instantânea e as definições de $P$
e $Q$, e pede $Q$. Ela envolve conceitos de \textbf{potência em circuitos de
corrente alternada}:

\begin{enumerate}
  \item \passo{Potência instantânea.} O produto de uma tensão e uma corrente
        senoidais de mesma frequência é uma senoide de frequência dupla
        somada a um valor constante. Escrever $p(t)$ explicitamente mostra
        quais características do gráfico correspondem a $P$ e a $VI$.
  \item \passo{Potências ativa, reativa e aparente.} $P$, $Q$ e a potência
        aparente $S = VI$ formam o triângulo de potências. A relação $S^2 =
        P^2 + Q^2$ dá $Q$ a partir de $P$ e $S$.
  \item \passo{Leitura de gráfico.} O valor médio e a amplitude de uma
        senoide deslocada do zero saem do seu máximo e do seu mínimo.
\end{enumerate}

\subsection{Revisão do conteúdo}

\subsubsection{Potência instantânea com tensão e corrente senoidais}

A potência instantânea entregue pelo gerador é o produto da tensão pela
corrente \cite{stevenson1986elementos}:
\begin{equation}
  p(t) = v(t)\,i(t) = 2VI\,\sen(\omega t)\,\sen(\omega t - \theta).
  \label{eq:q45-p-produto}
\end{equation}
Para transformar o produto de senos numa soma, parte-se das fórmulas do
cosseno da diferença e da soma:
\[
  \cos(a - b) = \cos a\cos b + \sen a\sen b,
  \qquad
  \cos(a + b) = \cos a\cos b - \sen a\sen b .
\]
Subtraindo a segunda da primeira, os termos $\cos a\cos b$ se cancelam:
\begin{equation}
  2\sen a\,\sen b = \cos(a - b) - \cos(a + b).
  \label{eq:q45-identidade}
\end{equation}
Com $a = \omega t$ e $b = \omega t - \theta$, tem-se $a - b = \theta$ e $a +
b = 2\omega t - \theta$. Substituindo na Equação~\eqref{eq:q45-p-produto}:
\begin{conceito}[title=Potência instantânea]
  \begin{equation}
    p(t) = VI\cos\theta - VI\cos(2\omega t - \theta)
    \label{eq:q45-p}
  \end{equation}
\end{conceito}

\paragraph{Valor médio.} Ao longo de um período $T = 2\pi/\omega$ da tensão,
o termo $\cos(2\omega t - \theta)$ completa exatamente dois ciclos, e sua
integral é nula:
\[
  \int_0^T \cos(2\omega t - \theta)\,dt
  = \left[\frac{\sen(2\omega t - \theta)}{2\omega}\right]_0^T
  = \frac{\sen(4\pi - \theta) - \sen(-\theta)}{2\omega}
  = \frac{-\sen\theta + \sen\theta}{2\omega} = 0 .
\]
Logo o valor médio de $p(t)$ é só o termo constante:
\[
  \frac{1}{T}\int_0^T p(t)\,dt = VI\cos\theta - 0 = VI\cos\theta = P,
\]
exatamente a definição de potência ativa dada no enunciado.

\paragraph{Amplitude, máximo e mínimo.} O termo oscilante tem amplitude
$VI$, que é a \emph{potência aparente} $S = VI$. Como o cosseno varia entre
$-1$ e $+1$:
\begin{equation}
  p_{\max} = VI\cos\theta + VI = P + S,
  \qquad
  p_{\min} = VI\cos\theta - VI = P - S .
  \label{eq:q45-max-min}
\end{equation}
Sempre que $\theta \neq 0$, tem-se $P < S$ e portanto $p_{\min} < 0$: em
parte do ciclo a potência instantânea é \emph{negativa}, isto é, a energia
volta da carga para o gerador.

\paragraph{Decomposição em parcela ativa e parcela reativa.} Expandindo o
cosseno da diferença,
\[
  \cos(2\omega t - \theta) = \cos 2\omega t\,\cos\theta
                           + \sen 2\omega t\,\sen\theta ,
\]
e substituindo na Equação~\eqref{eq:q45-p}:
\begin{equation}
  p(t) = VI\cos\theta\,(1 - \cos 2\omega t) - VI\sen\theta\,\sen 2\omega t
       = P\,(1 - \cos 2\omega t) - Q\,\sen 2\omega t .
  \label{eq:q45-decomposicao}
\end{equation}
\begin{itemize}
  \item A primeira parcela nunca é negativa: varia entre $0$ e $2P$ e tem
        média $P$. É a energia que vai do gerador para a carga e é
        consumida.
  \item A segunda parcela tem média nula e amplitude $Q$. É a energia que
        vai e volta a cada meio ciclo, armazenada e devolvida pelos campos
        magnéticos (indutores) e elétricos (capacitores). Ela não realiza
        trabalho líquido, mas circula pela rede. Esse é o significado físico
        da potência reativa.
\end{itemize}

\subsubsection{Potências ativa, reativa e aparente}

As três potências e suas unidades são:
\[
  S = VI~\text{(VA)}, \qquad
  P = S\cos\theta~\text{(W)}, \qquad
  Q = S\sen\theta~\text{(var)}.
\]
Elevando $P$ e $Q$ ao quadrado e somando, com $\cos^2\theta + \sen^2\theta =
1$:
\begin{conceito}[title=Triângulo de potências]
  \begin{equation}
    P^2 + Q^2 = S^2\left(\cos^2\theta + \sen^2\theta\right) = S^2
    \quad\Longrightarrow\quad
    Q = \pm\sqrt{S^2 - P^2}
    \label{eq:q45-triangulo}
  \end{equation}
\end{conceito}
$P$, $Q$ e $S$ são os lados de um triângulo retângulo
(Figura~\ref{fig:q45-triangulo}), e $\cos\theta = P/S$ é o \emph{fator de
potência}.

\begin{figure}[H]
  \centering
  \trianguloPotencias
  \caption{Triângulo de potências.}
  \label{fig:q45-triangulo}
\end{figure}

\paragraph{Sinal de $Q$.} Em $i(t) = \sqrt{2}\,I\sen(\omega t - \theta)$, o
ângulo $\theta$ é o atraso da corrente em relação à tensão. Com $\theta > 0$
a corrente está atrasada: a carga é indutiva e $Q > 0$. Com $\theta < 0$ a
corrente está adiantada: a carga é capacitiva e $Q < 0$. Por isso a raiz da
Equação~\eqref{eq:q45-triangulo} dá só o módulo de $Q$, e o sinal precisa
vir de outra informação.

\paragraph{Ligação com fasores.} Com $V = V\angle 0^\circ$ e $I =
I\angle{-\theta}$, a potência complexa da Equação~\eqref{eq:S} (Questão 24)
fica
\[
  V\,I^{*} = VI\angle\theta = VI\cos\theta + j\,VI\sen\theta = P + jQ :
\]
$P$ e $Q$ são as partes real e imaginária da potência complexa, e $S = VI$
é o seu módulo.

\subsubsection{Média e amplitude de uma senoide deslocada}

Uma curva da forma $y(t) = A - B\cos(\cdot)$, com $B > 0$, oscila entre
$y_{\max} = A + B$ e $y_{\min} = A - B$. Somando e subtraindo essas duas
relações:
\begin{equation}
  A = \frac{y_{\max} + y_{\min}}{2}, \qquad B = \frac{y_{\max} - y_{\min}}{2}.
  \label{eq:q45-media-amplitude}
\end{equation}
O valor médio ($A$) fica no meio entre o máximo e o mínimo, e a amplitude
($B$) é metade da distância entre eles.

\subsection{Resolução}

\subsubsection{Leitura do enunciado e do gráfico}

\begin{itemize}
  \item \passo{Intervalo de tempo.} O gráfico vai de $0$ a $0{,}02$~s, que
        é o intervalo $2\pi/\omega$ (um período da tensão). Então $\omega =
        2\pi/0{,}02 = 100\pi \approx 314$~rad/s, ou $f = 50$~Hz. Nesse
        intervalo, $p(t)$ completa \emph{dois} ciclos (dois máximos e dois
        mínimos), como esperado para uma oscilação em $2\omega$.
  \item \passo{Valores extremos.} Os picos ficam entre as linhas de 1400 e
        1600, e os vales entre $-400$ e $-600$:
        $p_{\max} \approx 1500$~W e $p_{\min} \approx -500$~W. A questão
        pede um valor \enquote{aproximado}, então essa precisão de leitura
        basta.
  \item \passo{O que não é dado.} $V$, $I$ e $\theta$ não aparecem
        separadamente; só a curva.
\end{itemize}

\subsubsection{Estratégia}

Comparando a Equação~\eqref{eq:q45-p}, $p(t) = P - S\cos(2\omega t -
\theta)$, com $y = A - B\cos(\cdot)$, o valor médio é $A = P$ e a amplitude
é $B = S$. A Equação~\eqref{eq:q45-media-amplitude} dá $P$ e $S$ diretamente
do máximo e do mínimo lidos no gráfico. Depois, $\cos\theta = P/S$ dá
$\theta$, e $Q = S\sen\theta$. Também é possível usar diretamente $Q =
\sqrt{S^2 - P^2}$, conferindo o sinal pela forma da curva.

\subsubsection{Passo 1: potência ativa}

\[
  P = \frac{p_{\max} + p_{\min}}{2}
    = \frac{1500 + (-500)}{2}
    = \frac{1000}{2} = 500~\text{W}.
\]

\subsubsection{Passo 2: potência aparente}

\[
  S = VI = \frac{p_{\max} - p_{\min}}{2}
    = \frac{1500 - (-500)}{2}
    = \frac{2000}{2} = 1000~\text{VA}.
\]

\subsubsection{Passo 3: ângulo e potência reativa}

O fator de potência e o ângulo:
\[
  \cos\theta = \frac{P}{S} = \frac{500}{1000} = 0{,}5
  \quad\Longrightarrow\quad
  \theta = \arccos 0{,}5 = 60^\circ ,
\]
com $\theta$ positivo, como se confere adiante. A potência reativa:
\[
  Q = S\sen\theta = 1000 \times \sen 60^\circ
    = 1000 \times \frac{\sqrt{3}}{2}
    = 500\sqrt{3} \approx 866~\text{var}.
\]
Pelo triângulo de potências (Equação~\eqref{eq:q45-triangulo}), o mesmo
valor:
\[
  Q = \sqrt{S^2 - P^2} = \sqrt{1000^2 - 500^2}
    = \sqrt{1\,000\,000 - 250\,000} = \sqrt{750\,000} \approx 866~\text{var}.
\]
O valor mais próximo entre as alternativas é 870~var.

\begin{resposta}
  \centering
  Alternativa \textbf{(B)}:\quad $Q = 500\sqrt{3} \approx 866$~var $\approx
  870$~var
\end{resposta}

\subsubsection{Conferências}

\passo{Sinal de $\theta$.} Pela Equação~\eqref{eq:q45-p}, $p(0) = VI\cos\theta
- VI\cos(-\theta) = 0$, como no gráfico. A derivada em $t = 0$ é
\[
  p'(0) = 2\omega VI\sen(2\omega\cdot 0 - \theta) = -2\omega VI\sen\theta .
\]
No gráfico a curva sai de zero \emph{descendo}, ou seja, $p'(0) < 0$, e
portanto $\sen\theta > 0$. Então $\theta = +60^\circ$ (não $-60^\circ$): a
corrente está atrasada, a carga é indutiva e $Q = +866$~var.

\passo{Reconstrução da curva.} Com $P = 500$~W, $S = 1000$~VA, $\theta =
60^\circ$ e $\omega = 100\pi$~rad/s (ou seja, $2\omega t = 36\,000\,t$ em
graus), a Equação~\eqref{eq:q45-p} fica
\[
  p(t) = 500 - 1000\cos\bigl(2\omega t - 60^\circ\bigr).
\]
Os pontos notáveis batem com o gráfico:

\begin{table}[H]
  \centering
  \begin{tabular}{@{}cccc@{}}
    \toprule
    $t$ (s) & $2\omega t - 60^\circ$ & $p(t)$ (W) & No gráfico \\
    \midrule
    $0$                        & $-60^\circ$  & $500 - 500 = 0$       & começa em zero \\
    $1/600 \approx 0{,}0017$   & $0^\circ$    & $500 - 1000 = -500$   & primeiro mínimo \\
    $1/300 \approx 0{,}0033$   & $60^\circ$   & $500 - 500 = 0$       & cruza o zero subindo \\
    $1/150 \approx 0{,}0067$   & $180^\circ$  & $500 + 1000 = 1500$   & primeiro máximo \\
    $1/100 = 0{,}01$           & $300^\circ$  & $500 - 500 = 0$       & cruza o zero descendo \\
    $7/600 \approx 0{,}0117$   & $360^\circ$  & $500 - 1000 = -500$   & segundo mínimo \\
    \bottomrule
  \end{tabular}
  \caption{Pontos notáveis de $p(t) = 500 - 1000\cos(2\omega t - 60^\circ)$.}
  \label{tab:q45-pontos}
\end{table}

\passo{Decomposição.} Pela Equação~\eqref{eq:q45-decomposicao}, $p(t) =
500\,(1 - \cos 2\omega t) - 866\,\sen 2\omega t$. A
Figura~\ref{fig:q45-decomposicao} mostra que a soma das duas parcelas
reproduz a curva do enunciado. A parcela ativa oscila entre 0 e 1000~W, com
média 500~W; a parcela reativa oscila entre $-866$ e $+866$~W, com média nula.

\begin{figure}[H]
  \centering
  \decomposicaoQquarentacinco
  \caption{A curva $p(t)$ do enunciado e as suas parcelas ativa e reativa.}
  \label{fig:q45-decomposicao}
\end{figure}

\passo{Por que as outras alternativas estão erradas.}
\begin{itemize}
  \item (A) 1\,000 é a potência aparente $S = VI$, não a reativa.
  \item (C) 500 é a potência ativa $P$ (o valor médio da curva).
  \item (D) 290 seria obtido com o ângulo errado, por exemplo $Q =
        P\,\mathrm{tg}\,30^\circ \approx 289$.
  \item (E) É verdade que o gráfico não permite achar $V$ e $I$
        separadamente: qualquer par com $VI = 1000$ produz a mesma curva.
        Mas $P$ e $Q$ dependem só do produto $VI$ e do ângulo $\theta$, e os
        dois saem do gráfico. Os dados são suficientes.
\end{itemize}

\begin{observacao}[title=Interpretação física]
  Nos intervalos em que a curva fica abaixo de zero ($0 < t < 1/300$~s e
  $0{,}01 < t < 0{,}0133$~s), a potência instantânea é negativa: o gerador
  recebe energia de volta da carga. É a energia que estava armazenada no
  campo magnético da carga indutiva ($\theta > 0$) sendo devolvida. A
  potência reativa, $Q \approx 866$~var, mede a amplitude dessa troca. O
  fator de potência $\cos\theta = 0{,}5$ é baixo: para entregar 500~W, o
  gerador precisa fornecer 1000~VA.
\end{observacao}

% =====================================================================
\newpage
\section*{Resumo das respostas}
\addcontentsline{toc}{section}{Resumo das respostas}

\begin{table}[H]
  \centering
  \begin{tabular}{@{}ll>{\raggedright\arraybackslash}p{7.6cm}@{}}
    \toprule
    Questão & Resposta & Ideia central \\
    \midrule
    51 & $P' = 5/6 \approx 0{,}833$~p.u. &
      $P \propto 1/X_{eq}$; $X_{eq}$ vai de $0{,}5$ para $0{,}6$~pu \\
    24 & Alternativa (D) &
      $I_{32} = (S_{G3}/V_3)^{*}$ $\Rightarrow$ $V_2$ $\Rightarrow$
      $I_{12}$ $\Rightarrow$ $S_{12}$ \\
    45 & Alternativa (B) &
      média $= P = 500$~W, amplitude $= S = 1000$~VA
      $\Rightarrow$ $Q = \sqrt{S^2 - P^2} \approx 866$~var \\
    \bottomrule
  \end{tabular}
\end{table}

\printbibliography[heading=bibintoc,title={Referências}]

\end{document}
